\bibitem{CRSBS}
J.~Burke and G.~Stevenson, \emph{The derived category of a graded Gorenstein ring},
\href{https://arxiv.org/abs/1507.00830}{arXiv:1507.00830v1} (2015),
Proposition~4.18, Lemmas~6.1--6.2 and Theorem~6.4.
PDF pp.~12--14, 17--19を確認。
\bibitem{CRSTV}
B.~To\"en and M.~Vaqui\'e, \emph{Moduli of objects in dg-categories},
Ann. Sci. \'Ecole Norm. Sup. (4) \textbf{40} (2007), 387--444;
\href{https://arxiv.org/abs/math/0503269}{arXiv:math/0503269v5}, Lemma~2.8(2).
PDF pp.~21--22を確認。pseudo-perfectからperfectへの原理を用いる。
\bibitem{CRSPerry}
A.~Perry, \emph{Noncommutative homological projective duality},
Adv. Math. \textbf{350} (2019), 877--972, Lemma~4.15(3).
\href{https://www.sciencedirect.com/science/article/am/pii/S000187081930221X}{著者受理稿}。
admissible半直交成分とsmoothnessの同値を用いる。
\bibitem{CRSStevenson}
G.~Stevenson, \emph{Subcategories of singularity categories via tensor actions},
Compos. Math. \textbf{150} (2014), 229--272;
\href{https://arxiv.org/abs/1105.4698}{arXiv:1105.4698v3}, Theorem~6.13 (PDF p.~26).
今回再確認したのはaffine hypersurfaceの全thick部分圏の分類である。
